% GENERATED FILE. Edit canonical lessons, then run npm run generate:books.
% canonical-source-hash: fnv1a64:19cc464a93892e8b
% canonical-lessons: SA-C85-week, SA-C85-monday, SA-C85-tuesday, SA-C85-wednesday, SA-C85-thursday

\chapter{Days of the Week}
\label{ch:sa-days-of-the-week}
\hlchaptermodality{85}

\begin{chapteropening}
\textbf{By the end of this chapter:} \emph{I can name the week, and Monday to Thursday.}

Complete the last lesson of chapter 85: I can name the week, and Monday to Thursday.
\end{chapteropening}

\section[\texorpdfstring{\sk{सप्ताहः} (saptāhaḥ)}{सप्ताहः (saptāhaḥ)}]{\sk{सप्ताहः} (saptāhaḥ) — a week}

\label{lesson:SA-C85-week}



\begin{quote}
Before the new one: say the Sanskrit for a farmer, then the Sanskrit for a merchant.
\end{quote}

\subsection*{You'll want to know: \sk{सप्ताहः}}
\textbf{\sk{सप्ताहः}} (\emph{saptāhaḥ}) — \textquotedblleft{}a week\textquotedblright{}: \textbf{\sk{सप्त}}, \textquotedblleft{}seven\textquotedblright{}, with \textbf{\sk{अह}}, \textquotedblleft{}day\textquotedblright{}.

\begin{cousinweb}
\emph{Saptāha} is built from \emph{sapta}, \textquotedblleft{}seven\textquotedblright{}, and \emph{aha}, \textquotedblleft{}day\textquotedblright{} (MW, CDIAL 13161). In the Brāhmaṇas it names a seven-day rite; the plain sense \textquotedblleft{}week\textquotedblright{} is later.
\end{cousinweb}

\begin{grammarlens}[title={What you've built}]
Seven plus day.
\end{grammarlens}

\subsection*{Your turn}
\begin{itemize}
\raggedright
  \item \emph{Say it:} \emph{saptāhaḥ}
  \item \emph{Say it:} \emph{saptāhaḥ}, once more
  \item \emph{Say it:} say \emph{saptāhaḥ}
  \item \emph{Recall:} say the Sanskrit for a farmer, then the Sanskrit for a merchant, then say \emph{saptāhaḥ} again
\end{itemize}

\begin{tcolorbox}[breakable,colback=teal!4,colframe=teal!35!black,title={Before you move on}]
What does \sk{सप्ताहः} mean? (\textquotedblleft{}A week\textquotedblright{}.) Say it once more.
\end{tcolorbox}
\section[\texorpdfstring{\sk{सोमवासरः} (somavāsaraḥ)}{सोमवासरः (somavāsaraḥ)}]{\sk{सोमवासरः} (somavāsaraḥ) — Monday}

\label{lesson:SA-C85-monday}



\begin{quote}
Before the new one: say the Sanskrit for a merchant, then the Sanskrit for a week.
\end{quote}

\subsection*{You'll want to know: \sk{सोमवासरः}}
\textbf{\sk{सोमवासरः}} (\emph{somavāsaraḥ}) — \textquotedblleft{}Monday\textquotedblright{}: \textbf{\sk{सोम}}, \textquotedblleft{}the moon\textquotedblright{}, with \textbf{\sk{वासरः}}, \textquotedblleft{}day\textquotedblright{}.

\begin{cousinweb}
\emph{Somavāsara} is built from \emph{soma}, the pressed drink and the moon, and \emph{vāsara}, \textquotedblleft{}day\textquotedblright{}, a word MW derives from the root √vas, \textquotedblleft{}to shine\textquotedblright{}. MW records the compound as \textquotedblleft{}Monday\textquotedblright{}.
\end{cousinweb}

\begin{grammarlens}[title={What you've built}]
The moon's day.
\end{grammarlens}

\subsection*{Your turn}
\begin{itemize}
\raggedright
  \item \emph{Say it:} \emph{somavāsaraḥ}
  \item \emph{Say it:} \emph{somavāsaraḥ}, once more
  \item \emph{Say it:} say \emph{somavāsaraḥ}
  \item \emph{Recall:} say the Sanskrit for a merchant, then the Sanskrit for a week, then say \emph{somavāsaraḥ} again
\end{itemize}

\begin{tcolorbox}[breakable,colback=teal!4,colframe=teal!35!black,title={Before you move on}]
What does \sk{सोमवासरः} mean? (\textquotedblleft{}Monday\textquotedblright{}.) Say it once more.
\end{tcolorbox}
\section[\texorpdfstring{\sk{मंगलवासरः} (maṅgalavāsaraḥ)}{मंगलवासरः (maṅgalavāsaraḥ)}]{\sk{मंगलवासरः} (maṅgalavāsaraḥ) — Tuesday}

\label{lesson:SA-C85-tuesday}



\begin{quote}
Before the new one: say the Sanskrit for a week, then the Sanskrit for Monday.
\end{quote}

\subsection*{You'll want to know: \sk{मंगलवासरः}}
\textbf{\sk{मंगलवासरः}} (\emph{maṅgalavāsaraḥ}) — \textquotedblleft{}Tuesday\textquotedblright{}, the day of \textbf{\sk{मंगल}}, Mars.

\begin{cousinweb}
\emph{Maṅgalavāsara} is built from \emph{maṅgala}, \textquotedblleft{}welfare, auspiciousness\textquotedblright{}, the name given to the planet Mars, and \emph{vāsara}, \textquotedblleft{}day\textquotedblright{} (MW).
\end{cousinweb}

\begin{grammarlens}[title={What you've built}]
Mars's day.
\end{grammarlens}

\subsection*{Your turn}
\begin{itemize}
\raggedright
  \item \emph{Say it:} \emph{maṅgalavāsaraḥ}
  \item \emph{Say it:} \emph{maṅgalavāsaraḥ}, once more
  \item \emph{Say it:} say \emph{maṅgalavāsaraḥ}
  \item \emph{Recall:} say the Sanskrit for a week, then the Sanskrit for Monday, then say \emph{maṅgalavāsaraḥ} again
\end{itemize}

\begin{tcolorbox}[breakable,colback=teal!4,colframe=teal!35!black,title={Before you move on}]
What does \sk{मंगलवासरः} mean? (\textquotedblleft{}Tuesday\textquotedblright{}.) Say it once more.
\end{tcolorbox}
\section[\texorpdfstring{\sk{बुधवासरः} (budhavāsaraḥ)}{बुधवासरः (budhavāsaraḥ)}]{\sk{बुधवासरः} (budhavāsaraḥ) — Wednesday}

\label{lesson:SA-C85-wednesday}



\begin{quote}
Before the new one: say the Sanskrit for Monday, then the Sanskrit for Tuesday.
\end{quote}

\subsection*{You'll want to know: \sk{बुधवासरः}}
\textbf{\sk{बुधवासरः}} (\emph{budhavāsaraḥ}) — \textquotedblleft{}Wednesday\textquotedblright{}, the day of \textbf{\sk{बुध}}, Mercury.

\begin{cousinweb}
\emph{Budhavāsara} is built from \emph{budha}, \textquotedblleft{}awake, wise\textquotedblright{}, from the root √budh, \textquotedblleft{}to wake\textquotedblright{}, the name given to Mercury, and \emph{vāsara}, \textquotedblleft{}day\textquotedblright{} (MW). MW records the day as \emph{budhavāra}, with \emph{vāra}, \textquotedblleft{}turn\textquotedblright{}, in place of \emph{vāsara}.
\end{cousinweb}

\begin{grammarlens}[title={What you've built}]
Mercury's day.
\end{grammarlens}

\subsection*{Your turn}
\begin{itemize}
\raggedright
  \item \emph{Say it:} \emph{budhavāsaraḥ}
  \item \emph{Say it:} \emph{budhavāsaraḥ}, once more
  \item \emph{Say it:} say \emph{budhavāsaraḥ}
  \item \emph{Recall:} say the Sanskrit for Monday, then the Sanskrit for Tuesday, then say \emph{budhavāsaraḥ} again
\end{itemize}

\begin{tcolorbox}[breakable,colback=teal!4,colframe=teal!35!black,title={Before you move on}]
What does \sk{बुधवासरः} mean? (\textquotedblleft{}Wednesday\textquotedblright{}.) Say it once more.
\end{tcolorbox}
\section[\texorpdfstring{\sk{गुरुवासरः} (guruvāsaraḥ)}{गुरुवासरः (guruvāsaraḥ)}]{\sk{गुरुवासरः} (guruvāsaraḥ) — Thursday}

\label{lesson:SA-C85-thursday}



\begin{quote}
Before the new one: say the Sanskrit for Tuesday, then the Sanskrit for Wednesday.
\end{quote}

\subsection*{You'll want to know: \sk{गुरुवासरः}}
\textbf{\sk{गुरुवासरः}} (\emph{guruvāsaraḥ}) — \textquotedblleft{}Thursday\textquotedblright{}, the day of Jupiter, called \textbf{\sk{गुरुः}}, \textquotedblleft{}the teacher\textquotedblright{} — a word you know.

\begin{cousinweb}
\emph{Guruvāsara} is built from \emph{guru}, \textquotedblleft{}heavy, the teacher\textquotedblright{}, the name given to Jupiter, and \emph{vāsara}, \textquotedblleft{}day\textquotedblright{} (MW). MW records the day as \emph{guruvāra}, with \emph{vāra}, \textquotedblleft{}turn\textquotedblright{}, in place of \emph{vāsara}.
\end{cousinweb}

\begin{grammarlens}[title={What you've built}]
The teacher's day.
\end{grammarlens}

\subsection*{Your turn}
\begin{itemize}
\raggedright
  \item \emph{Say it:} \emph{guruvāsaraḥ}
  \item \emph{Say it:} \emph{guruvāsaraḥ}, once more
  \item \emph{Say it:} say \emph{guruvāsaraḥ}
  \item \emph{Recall:} say the Sanskrit for Tuesday, then the Sanskrit for Wednesday, then say \emph{guruvāsaraḥ} again
\end{itemize}

\begin{tcolorbox}[breakable,colback=teal!4,colframe=teal!35!black,title={Before you move on}]
What does \sk{गुरुवासरः} mean? (\textquotedblleft{}Thursday\textquotedblright{}.) Say it once more.
\end{tcolorbox}
